% torsion_realizations.tex — torsion subgroups with a geometrically simple
% genus-2 Jacobian realization over Q, with minimal-conductor examples and
% earliest references.  Generated by code/claude_gen_torsion_table.py.
\documentclass[11pt]{amsart}
\usepackage{array}
\usepackage{longtable}
\usepackage[margin=2cm]{geometry}
\usepackage[colorlinks=true,linkcolor=blue,citecolor=blue,urlcolor=blue]{hyperref}
\title{Exact torsion groups realized on geometrically simple\\
genus-2 Jacobians over $\mathbb{Q}$}
\date{July 2026}
\begin{document}
\maketitle
\noindent Groups are written by invariant factors: $[n_1,\ldots,n_r]$ denotes
$\mathbb{Z}/n_1\mathbb{Z}\times\cdots\times\mathbb{Z}/n_r\mathbb{Z}$ with
$n_1\mid\cdots\mid n_r$.  Each row records one genus-2 curve $C/\mathbb{Q}$ whose
Jacobian is geometrically simple and whose \emph{full} rational torsion group is
\emph{exactly} the displayed group, $\operatorname{Jac}(C)(\mathbb{Q})_{\rm
tors}\cong G$ (the first row is the trivial group, written $[\,]$).  The \emph{Source(s)} column, by
contrast, dates the \emph{group}, and the earliest source there may establish
only a rational point of the given order, a subgroup, or a positive-dimensional
family containing $G$ --- not necessarily an exact-torsion or simple example; the
companion document \texttt{torsion\_sources.tex} classifies every source
individually.  We display, for each group, the equation of the curve of smallest
conductor in the queried database snapshot (2026-07-22) with $\operatorname{Jac}$
geometrically simple and (where such a curve exists; see the endomorphism note)
generic endomorphisms.  Curves of small conductor lie in the production LMFDB;
those found only in the extended database of Booker--Sutherland cite
\cite{BookerSutherland} in the \emph{Source(s)} column.  For convenience each
equation is hyperlinked to its LMFDB home page: production curves by their
permanent label, and curves that are only in the extended database via a
alpha-site (\url{https://alpha.lmfdb.org}) equation-lookup URL, since alpha
labels are not permanent (the lookup matches any isomorphic model of the
curve).  The equations are the intended stable identifiers.

\emph{Source(s) column.} It gives the earliest construction of the torsion,
together with the earliest certification of geometric simplicity when these
differ; ``new'' means the exact simple realization first appears in the present
paper.  A dagger ($\dagger$) marks a source that exhibited the torsion but did
not establish geometric simplicity of a realizing curve; in daggered rows the
certification is supplied by the linked database entry or by the co-cited later
source.  Lepr\'evost's classical order-21, 25 and 27 curves were later shown
non-simple or geometrically split, so those rows cite the first proven-simple
realizations: the certified database curve \texttt{388.a.776.1} realizes $[21]$,
Nicholls \cite{Nicholls2018} realizes $[25]$, and for $[27]$ both Nicholls
\cite{Nicholls2018} and the certified database curve \texttt{604.a.9664.2}
(each distinct from Lepr\'evost's split curve \cite{Leprevost1995}).
Lepr\'evost's $[3,9]$ family is proven
simple over $\mathbb{Q}$ only.  Elkies' page exhibits a rational order-31
\emph{subgroup} with no rational generator (not a realization); the first
\emph{explicit} order-31 realization --- previously the smallest cyclic order
with no explicit realization, and the first on a genus-2 Jacobian --- is the
$[31]$ row: the curve attached to the modular form
\texttt{1830.2.a.q} from the modular-abelian-surfaces dataset \cite{CEHJMPV},
whose exact rational torsion ($\mathbb{Z}/31\mathbb{Z}$, with generator of
irrational Mumford support) and geometric simplicity are certified here.
Order 31 also appears (for $\mathrm{GL}_2$-type abelian varieties, without
curve equations) in the torsion conjectures of Alessandr\`i--Coppola
\cite{AC2026}, whose Theorem 4.2/4.3 proves an order-31 point exists on some
inexplicit member of this isogeny class --- the novelty here is the explicit
curve, its exact group, and its certificates; their $g=2$ lists include 37 as
well, but at the one
admissible level ($2190$, newform \texttt{2190.2.a.v}) no member of the class
is a Jacobian over $\mathbb{Q}$ --- the real multiplication field
$\mathbb{Q}(\sqrt3)$ has narrow class number $2$, so no member carries a
rational principal polarization.

\emph{Endomorphism refinement.} With two exceptions, every group here is realized
by a curve with \emph{generic} geometric endomorphisms,
$\operatorname{End}(\operatorname{Jac}(C)_{\overline{\mathbb{Q}}})=\mathbb{Z}$
(the displayed curve is the smallest-conductor known such witness).  For database
rows this is the LMFDB endomorphism certification (Costa--Mascot--Sijsling--Voight).
The exceptions, marked $^{\rm RM}$, are $[2,22]$ and $[31]$: no generic witness is
known for either.  For $[2,22]$ the extended database's only geometrically simple
curves have real multiplication ($\operatorname{End}_{\overline{\mathbb{Q}}}$ a
real quadratic order), so the displayed curve is geometrically simple but not
generic.  For $[31]$ the only known witness is the displayed curve attached to
the newform \texttt{1830.2.a.q} (real multiplication by $\mathbb{Z}[\sqrt2]$,
conductor $1830^2 = 3348900$, outside the extended database's conductor range);
its $31$-torsion arises from an Eisenstein congruence at the level prime
$61$ ($31 \mid 61+1$).  For
$[2,2,14]$ the extended database likewise contains only RM witnesses; the
displayed curve --- of conductor
$2^2 3^2 5^2 7^2 13^2\cdot 19\cdot 23\cdot 37\cdot 317 = 38200383011700$, the
smallest of ten generic witnesses found here by a three-root refinement of a
contact-7 construction --- is new to this paper.
For each displayed non-database
equation with generic endomorphisms (the rows marked ``new'' without an
$^{\rm RM}$ flag, and the $[2,2,2,10]$ family member) we give a
self-contained certificate in \texttt{data/claude\_endz\_certificates.txt}: a
good prime whose $L$-polynomial is irreducible with all Frobenius-root powers of
degree $4$ (geometric simplicity, by the root-power criterion), a second good
prime whose irreducible $L$-polynomial has splitting field linearly disjoint
from the first (so the centre of the geometric endomorphism algebra is
$\mathbb{Q}$, excluding real and complex multiplication; cf.\
\cite{CLV2021,Zywina2022}), and the observation that $\#\operatorname{Jac}(C)(
\mathbb{Q})_{\rm tors}>18$, which excludes quaternionic multiplication since an
(O-)PQM genus-2 Jacobian over $\mathbb{Q}$ has at most $16$ (resp.\ $18$) rational
torsion points \cite{LSSV2024}.  The RM-flagged non-database row $[31]$ cannot
carry such an $\operatorname{End}=\mathbb{Z}$ certificate (its geometric
endomorphism ring is $\mathbb{Z}[\sqrt2]$); its dedicated certificate --- exact
torsion invariants, the $D_4$/root-power geometric-simplicity witness at
$p=11$, and the real-multiplication identification (constant Frobenius
discriminant core $2$ at all $164$ good primes $p \le 1000$, Hecke-trace match
to \texttt{1830.2.a.q}) --- is in \texttt{data/claude\_z31\_rm\_certificate.txt}.
The one caveat is that the earliest cited simple
realization is occasionally itself non-generic --- $\ddagger$ marks $[5]$, whose
first simple realizations are the RM Jacobian $J_0(31)$ and the CM curve
$y^2-y=x^5$; the displayed witness in every such row is nonetheless generic.
\medskip

\begingroup\small
\renewcommand{\arraystretch}{1.25}
\begin{longtable}{|>{\raggedright\arraybackslash}p{1.7cm}|>{\raggedright\arraybackslash}p{0.70\textwidth}|>{\raggedright\arraybackslash}p{2.1cm}|}
\hline
Group & Curve & Source(s)\\ \hline
\endfirsthead
\hline
Group & Curve & Source(s)\\ \hline
\endhead
$[\,]$ & \href{https://alpha.lmfdb.org/Genus2Curve/Q/?jump=[[-14580,51520,-63684,31054,-3823,-637,-18],[1]]}{$y^2 + y = -18x^{6}-637x^{5}-3823x^{4}+31054x^{3}-63684x^{2}+51520x-14580$} & \cite{BSSVY}\\ \hline
$[2]$ & \href{https://www.lmfdb.org/Genus2Curve/Q/295/a/295/2}{$y^2 + (x^{2}+x+1)y = x^{5}-40x^{3}+22x^{2}+389x-608$} & \cite{BSSVY}\\ \hline
$[3]$ & \href{https://www.lmfdb.org/Genus2Curve/Q/997/b/997/1}{$y^2 + y = x^{5}-2x^{4}+2x^{3}-x^{2}$} & \cite{BSSVY}\\ \hline
$[4]$ & \href{https://www.lmfdb.org/Genus2Curve/Q/1070/a/2140/1}{$y^2 + (x^{3}+1)y = x^{3}-x$} & \cite{BSSVY}\\ \hline
$[5]$ & \href{https://www.lmfdb.org/Genus2Curve/Q/277/a/277/2}{$y^2 + y = x^{5}-9x^{4}+14x^{3}-19x^{2}+11x-6$} & \cite{Ogg1973}$^\dagger$, \cite{BoxallGrant2000}$^\ddagger$\\ \hline
$[6]$ & \href{https://www.lmfdb.org/Genus2Curve/Q/1038/a/1038/1}{$y^2 + (x^{2}+x)y = x^{5}-12x^{4}+26x^{3}+46x^{2}+21x+3$} & \cite{Flynn1991}$^\dagger$\\ \hline
$[7]$ & \href{https://www.lmfdb.org/Genus2Curve/Q/461/a/461/1}{$y^2 + x^{3}y = x^{5}-3x^{3}+3x-2$} & \cite{Ogg1973}$^\dagger$, \cite{LPS2004}\\ \hline
$[8]$ & \href{https://www.lmfdb.org/Genus2Curve/Q/464/a/464/1}{$y^2 + (x+1)y = -x^{6}-2x^{5}-2x^{4}-x^{3}$} & \cite{BSSVY}\\ \hline
$[9]$ & \href{https://www.lmfdb.org/Genus2Curve/Q/713/b/713/1}{$y^2 + (x^{3}+x+1)y = -x^{4}$} & \cite{Flynn1991}$^\dagger$\\ \hline
$[10]$ & \href{https://www.lmfdb.org/Genus2Curve/Q/389/a/389/1}{$y^2 + (x^{3}+x)y = x^{5}-2x^{4}-8x^{3}+16x+7$} & \cite{Flynn1991}$^\dagger$\\ \hline
$[11]$ & \href{https://www.lmfdb.org/Genus2Curve/Q/353/a/353/1}{$y^2 + (x^{3}+x+1)y = x^{2}$} & \cite{Ogg1973}$^\dagger$, \cite{BLP2009}\\ \hline
$[12]$ & \href{https://www.lmfdb.org/Genus2Curve/Q/762/a/3048/1}{$y^2 + (x^{3}+x^{2}+x)y = x^{2}+x+1$} & \cite{BSSVY}\\ \hline
$[13]$ & \href{https://www.lmfdb.org/Genus2Curve/Q/349/a/349/1}{$y^2 + (x^{3}+x^{2}+x+1)y = -x^{3}-x^{2}$} & \cite{Flynn1990,Leprevost1991a}$^\dagger$, \cite{PZP2013}\\ \hline
$[14]$ & \href{https://www.lmfdb.org/Genus2Curve/Q/249/a/249/1}{$y^2 + (x^{3}+1)y = x^{2}+x$} & \cite{PP2012,PZP2013}\\ \hline
$[15]$ & \href{https://www.lmfdb.org/Genus2Curve/Q/277/a/277/1}{$y^2 + (x^{3}+x^{2}+x+1)y = -x^{2}-x$} & \cite{Leprevost1991b}$^\dagger$\\ \hline
$[16]$ & \href{https://www.lmfdb.org/Genus2Curve/Q/830/a/6640/1}{$y^2 + (x^{3}+1)y = -x^{5}+x^{4}-2x^{2}+x+1$} & \cite{BSSVY}\\ \hline
$[17]$ & \href{https://www.lmfdb.org/Genus2Curve/Q/1996/b/510976/1}{$y^2 + (x^{3}+x^{2}+x)y = x^{3}+x^{2}+3x+1$} & \cite{Leprevost1991b}$^\dagger$, \cite{PZP2013}\\ \hline
$[18]$ & \href{https://www.lmfdb.org/Genus2Curve/Q/1180/a/18880/1}{$y^2 + (x^{3}+1)y = -2x^{4}+4x^{2}+2x$} & \cite{PP2012,PZP2013}\\ \hline
$[19]$ & \href{https://alpha.lmfdb.org/Genus2Curve/Q/?jump=[[-2,5,-3,-5,1,1,1],[1,1,1]]}{$y^2 + (x^{2}+x+1)y = x^{6}+x^{5}+x^{4}-5x^{3}-3x^{2}+5x-2$} & \cite{Leprevost1991b}$^\dagger$, \cite{PZP2013}\\ \hline
$[20]$ & \href{https://www.lmfdb.org/Genus2Curve/Q/394/a/3152/1}{$y^2 + (x+1)y = -x^{5}$} & \cite{Leprevost1997}$^\dagger$\\ \hline
$[21]$ & \href{https://www.lmfdb.org/Genus2Curve/Q/388/a/776/1}{$y^2 + (x^{3}+x+1)y = -x^{4}+2x^{2}+x$} & \cite{Leprevost1991b}$^\dagger$, \cite{Nicholls2018}\\ \hline
$[22]$ & \href{https://www.lmfdb.org/Genus2Curve/Q/1192/a/19072/1}{$y^2 + (x^{3}+x)y = x^{3}-2x^{2}-x+1$} & \cite{Leprevost1993}$^\dagger$, \cite{Leprevost1995}\\ \hline
$[23]$ & \href{https://alpha.lmfdb.org/Genus2Curve/Q/?jump=[[1,-2,2,1,-2,-1,1],[0,0,1]]}{$y^2 + x^{2}y = x^{6}-x^{5}-2x^{4}+x^{3}+2x^{2}-2x+1$} & \cite{Ogawa1994}$^\dagger$, \cite{Leprevost1995}\\ \hline
$[24]$ & \href{https://www.lmfdb.org/Genus2Curve/Q/1908/a/183168/1}{$y^2 + (x^{3}+1)y = 2x^{4}+3x^{3}+4x^{2}+2x$} & \cite{Leprevost1993}$^\dagger$, \cite{Leprevost1995}\\ \hline
$[25]$ & \href{https://alpha.lmfdb.org/Genus2Curve/Q/?jump=[[2,-23,41,75,25,-9],[1,1]]}{$y^2 + (x+1)y = -9x^{5}+25x^{4}+75x^{3}+41x^{2}-23x+2$} & \cite{Leprevost1993}$^\dagger$, \cite{Nicholls2018}\\ \hline
$[26]$ & \href{https://alpha.lmfdb.org/Genus2Curve/Q/?jump=[[0,0,-3,-3,-3,-12],[1,1]]}{$y^2 + (x+1)y = -12x^{5}-3x^{4}-3x^{3}-3x^{2}$} & \cite{Leprevost1993}$^\dagger$, \cite{Leprevost1995}\\ \hline
$[27]$ & \href{https://www.lmfdb.org/Genus2Curve/Q/604/a/9664/2}{$y^2 + (x^{3}+1)y = -x^{4}+x^{3}+x^{2}-x$} & \cite{Howe2015}$^\dagger$, \cite{Nicholls2018}\\ \hline
$[28]$ & \href{https://www.lmfdb.org/Genus2Curve/Q/249/a/6723/1}{$y^2 + (x^{3}+1)y = -x^{5}+x^{3}+x^{2}+3x+2$} & \cite{PP2012,PZP2013}\\ \hline
$[29]$ & \href{https://www.lmfdb.org/Genus2Curve/Q/976/a/999424/1}{$y^2 + (x+1)y = x^{6}-2x^{5}+2x^{3}-x^{2}$} & \cite{Leprevost1993}$^\dagger$, \cite{Leprevost1995}\\ \hline
$[30]$ & \href{https://alpha.lmfdb.org/Genus2Curve/Q/?jump=[[0,-1,16,-12,-10,4,2],[0,0,1,1]]}{$y^2 + (x^{3}+x^{2})y = 2x^{6}+4x^{5}-10x^{4}-12x^{3}+16x^{2}-x$} & \cite{PP2015}\\ \hline
$[31]$ & $y^2 + (-x^{2}-x)y = -839x^{6}+2841x^{5}-4587x^{4}+4300x^{3}-2466x^{2}+816x-126$ & \cite{CEHJMPV}$^\dagger$, new$^{\rm RM}$\\ \hline
$[32]$ & \href{https://alpha.lmfdb.org/Genus2Curve/Q/?jump=[[-47,-49,49,45,1,-1,1],[0,1,1]]}{$y^2 + (x^{2}+x)y = x^{6}-x^{5}+x^{4}+45x^{3}+49x^{2}-49x-47$} & \cite{Elkies2002}\\ \hline
$[33]$ & \href{https://alpha.lmfdb.org/Genus2Curve/Q/?jump=[[2,13,21,7,9,-7,1],[1,1,1]]}{$y^2 + (x^{2}+x+1)y = x^{6}-7x^{5}+9x^{4}+7x^{3}+21x^{2}+13x+2$} & \cite{PZP2013}\\ \hline
$[34]$ & \href{https://alpha.lmfdb.org/Genus2Curve/Q/?jump=[[15,12,7,3,0,-2],[0,0,1,1]]}{$y^2 + (x^{3}+x^{2})y = -2x^{5}+3x^{3}+7x^{2}+12x+15$} & \cite{Elkies2002}\\ \hline
$[36]$ & \href{https://alpha.lmfdb.org/Genus2Curve/Q/?jump=[[9,6,-14,13,30,3],[0,0,1]]}{$y^2 + x^{2}y = 3x^{5}+30x^{4}+13x^{3}-14x^{2}+6x+9$} & \cite{PP2015}\\ \hline
$[39]$ & \href{https://www.lmfdb.org/Genus2Curve/Q/1116/a/214272/1}{$y^2 + (x^{3}+1)y = x^{4}+2x^{3}+x^{2}-x$} & \cite{Elkies2002}\\ \hline
$[40]$ & \href{https://alpha.lmfdb.org/Genus2Curve/Q/?jump=[[0,0,-2,-2,4,3],[1,1]]}{$y^2 + (x+1)y = 3x^{5}+4x^{4}-2x^{3}-2x^{2}$} & \cite{Elkies2002}\\ \hline
$[2,\,2]$ & \href{https://www.lmfdb.org/Genus2Curve/Q/464/a/29696/2}{$y^2 + xy = 4x^{5}+33x^{4}+72x^{3}+16x^{2}+x$} & \cite{BSSVY}\\ \hline
$[2,\,4]$ & \href{https://www.lmfdb.org/Genus2Curve/Q/997/a/997/1}{$y^2 + xy = x^{5}-8x^{4}+16x^{3}-x$} & \cite{BSSVY}\\ \hline
$[2,\,6]$ & \href{https://www.lmfdb.org/Genus2Curve/Q/704/a/45056/1}{$y^2 + y = 4x^{5}+4x^{4}-x^{3}-2x^{2}$} & \cite{BSSVY}\\ \hline
$[2,\,8]$ & \href{https://www.lmfdb.org/Genus2Curve/Q/464/a/29696/1}{$y^2 + (x+1)y = 8x^{5}+3x^{4}-4x^{3}-2x^{2}$} & \cite{BSSVY}\\ \hline
$[2,\,10]$ & \href{https://www.lmfdb.org/Genus2Curve/Q/555/a/8325/1}{$y^2 + (x+1)y = 3x^{5}-2x^{4}-4x^{3}+x^{2}+x$} & \cite{BSSVY}\\ \hline
$[2,\,12]$ & \href{https://www.lmfdb.org/Genus2Curve/Q/762/a/82296/1}{$y^2 + (x^{2}+x)y = x^{5}-8x^{4}+14x^{3}+2x^{2}-x$} & \cite{BSSVY}\\ \hline
$[2,\,14]$ & \href{https://www.lmfdb.org/Genus2Curve/Q/1416/b/135936/1}{$y^2 + (x^{3}+x)y = -2x^{4}-x^{3}+x+1$} & \cite{BSSVY}\\ \hline
$[2,\,16]$ & \href{https://alpha.lmfdb.org/Genus2Curve/Q/?jump=[[8,4,-3,1,-1,-1],[0,0,1,1]]}{$y^2 + (x^{3}+x^{2})y = -x^{5}-x^{4}+x^{3}-3x^{2}+4x+8$} & \cite{BookerSutherland}\\ \hline
$[2,\,18]$ & \href{https://alpha.lmfdb.org/Genus2Curve/Q/?jump=[[7,2,1,3,-5,-1,1],[0,1,1]]}{$y^2 + (x^{2}+x)y = x^{6}-x^{5}-5x^{4}+3x^{3}+x^{2}+2x+7$} & \cite{BookerSutherland}\\ \hline
$[2,\,20]$ & \href{https://alpha.lmfdb.org/Genus2Curve/Q/?jump=[[0,-180,-251,185,28,-30,4],[0,0,1]]}{$y^2 + x^{2}y = 4x^{6}-30x^{5}+28x^{4}+185x^{3}-251x^{2}-180x$} & \cite{BookerSutherland}\\ \hline
$[2,\,22]$ & \href{https://alpha.lmfdb.org/Genus2Curve/Q/?jump=[[0,-6,12,-5,9,-3,1],[0,1,1]]}{$y^2 + (x^{2}+x)y = x^{6}-3x^{5}+9x^{4}-5x^{3}+12x^{2}-6x$} & \cite{BookerSutherland}$^{\rm RM}$\\ \hline
$[2,\,26]$ & \href{https://alpha.lmfdb.org/Genus2Curve/Q/?jump=[[0,0,2,3,-12,-3,9],[1,0,1]]}{$y^2 + (x^{2}+1)y = 9x^{6}-3x^{5}-12x^{4}+3x^{3}+2x^{2}$} & \cite{BookerSutherland}\\ \hline
$[2,\,28]$ & \href{https://alpha.lmfdb.org/Genus2Curve/Q/?jump=[[0,-3,4,34,-42,-3,9],[1,0,1]]}{$y^2 + (x^{2}+1)y = 9x^{6}-3x^{5}-42x^{4}+34x^{3}+4x^{2}-3x$} & \cite{BookerSutherland}\\ \hline
$[3,\,3]$ & \href{https://alpha.lmfdb.org/Genus2Curve/Q/?jump=[[3,3,7,2,4,0,1],[1]]}{$y^2 + y = x^{6}+4x^{4}+2x^{3}+7x^{2}+3x+3$} & \cite{BFT2014}\\ \hline
$[3,\,6]$ & \href{https://alpha.lmfdb.org/Genus2Curve/Q/?jump=[[21,-2,7,-5,0,-1],[1,1,1,1]]}{$y^2 + (x^{3}+x^{2}+x+1)y = -x^{5}-5x^{3}+7x^{2}-2x+21$} & \cite{BookerSutherland}\\ \hline
$[3,\,9]$ & \href{https://alpha.lmfdb.org/Genus2Curve/Q/?jump=[[326,-75,-230,-113,40,75,16],[1,0,0,1]]}{$y^2 + (x^{3}+1)y = 16x^{6}+75x^{5}+40x^{4}-113x^{3}-230x^{2}-75x+326$} & \cite{Leprevost1995}$^\dagger$, \cite{BookerSutherland}\\ \hline
$[4,\,4]$ & \href{https://alpha.lmfdb.org/Genus2Curve/Q/?jump=[[20493,6534,10832,2601,1823,242,99],[0,1,1]]}{$y^2 + (x^{2}+x)y = 99x^{6}+242x^{5}+1823x^{4}+2601x^{3}+10832x^{2}+6534x+20493$} & \cite{BookerSutherland}\\ \hline
$[4,\,8]$ & \href{https://alpha.lmfdb.org/Genus2Curve/Q/?jump=[[164,682,450,-435,272,-81],[0,1,1]]}{$y^2 + (x^{2}+x)y = -81x^{5}+272x^{4}-435x^{3}+450x^{2}+682x+164$} & \cite{BookerSutherland}\\ \hline
$[6,\,6]$ & $y^2 = x(39x^2-69x+125)(48x^2+8x+39)$ & new\\ \hline
$[2,\,2,\,2]$ & \href{https://www.lmfdb.org/Genus2Curve/Q/2600/a/338000/1}{$y^2 + xy = 10x^{5}+8x^{4}-5x^{3}-3x^{2}+x$} & \cite{BSSVY}\\ \hline
$[2,\,2,\,4]$ & \href{https://www.lmfdb.org/Genus2Curve/Q/3978/a/930852/1}{$y^2 + (x^{2}+x)y = x^{5}+3x^{4}-3x^{3}-8x^{2}+6x$} & \cite{BSSVY}\\ \hline
$[2,\,2,\,6]$ & \href{https://www.lmfdb.org/Genus2Curve/Q/816/a/39168/1}{$y^2 + (x^{2}+1)y = 3x^{5}-4x^{3}-x^{2}+x$} & \cite{BSSVY}\\ \hline
$[2,\,2,\,8]$ & \href{https://alpha.lmfdb.org/Genus2Curve/Q/?jump=[[0,-1,-6,-3,20,9],[0,1,1]]}{$y^2 + (x^{2}+x)y = 9x^{5}+20x^{4}-3x^{3}-6x^{2}-x$} & \cite{BookerSutherland}\\ \hline
$[2,\,2,\,10]$ & \href{https://alpha.lmfdb.org/Genus2Curve/Q/?jump=[[0,0,-7,12,6,-12],[1,0,1]]}{$y^2 + (x^{2}+1)y = -12x^{5}+6x^{4}+12x^{3}-7x^{2}$} & \cite{BookerSutherland}\\ \hline
$[2,\,2,\,12]$ & \href{https://alpha.lmfdb.org/Genus2Curve/Q/?jump=[[0,-5,-5,16,8,-7,1],[0,1,1]]}{$y^2 + (x^{2}+x)y = x^{6}-7x^{5}+8x^{4}+16x^{3}-5x^{2}-5x$} & \cite{BookerSutherland}\\ \hline
$[2,\,2,\,14]$ & $y^2 = (x+1)(x+9)(2x-115)(6x+55)(3x^2-220x+2652)$ & \cite{BookerSutherland}, new\\ \hline
$[2,\,2,\,20]$ & $y^2 = -(x-1)(6x+1)(2x+7)(6217x+1008)(21x^2-161x+144)$ & new\\ \hline
$[2,\,4,\,4]$ & \href{https://alpha.lmfdb.org/Genus2Curve/Q/?jump=[[0,180,-210,-1850,2204,-116],[0,1,1]]}{$y^2 + (x^{2}+x)y = -116x^{5}+2204x^{4}-1850x^{3}-210x^{2}+180x$} & \cite{BookerSutherland}\\ \hline
$[2,\,4,\,8]$ & $y^2 = x(3x-5)(5x+333)(16x^2+23x+576)$ & new\\ \hline
$[2,\,2,\,2,\,2]$ & \href{https://alpha.lmfdb.org/Genus2Curve/Q/?jump=[[0,-1,-4,8,23,-27],[0,1,1]]}{$y^2 + (x^{2}+x)y = -27x^{5}+23x^{4}+8x^{3}-4x^{2}-x$} & \cite{BookerSutherland}\\ \hline
$[2,\,2,\,2,\,4]$ & \href{https://alpha.lmfdb.org/Genus2Curve/Q/?jump=[[0,1,-6,-16,146,-225],[0,1,1]]}{$y^2 + (x^{2}+x)y = -225x^{5}+146x^{4}-16x^{3}-6x^{2}+x$} & \cite{BookerSutherland}\\ \hline
$[2,\,2,\,2,\,6]$ & \href{https://alpha.lmfdb.org/Genus2Curve/Q/?jump=[[2,-12,11,18,-14,-6],[1,0,1]]}{$y^2 + (x^{2}+1)y = -6x^{5}-14x^{4}+18x^{3}+11x^{2}-12x+2$} & \cite{BookerSutherland}\\ \hline
$[2,\,2,\,2,\,8]$ & $y^2 = x(x+1)(x+55^2)(x+99^2)(x+125^2)$ & new\\ \hline
$[2,\,2,\,2,\,10]$ & $y^2 = x(x+1)(x-1)(3x-7)(8x-13)(24x+25)$ & \cite{Elkies2024}\\ \hline
$[2,\,2,\,2,\,12]$ & $y^2 = x(x-120^2)(x-143^2)(x-266^2)(x-218^2)(x-241^2)$ & new\\ \hline
$[2,\,2,\,4,\,4]$ & $y^2 = x(x+36^2)(x+57^2)(x+64^2)(x+132^2)$ & new\\ \hline
\end{longtable}
\endgroup

\begin{thebibliography}{99}
\bibitem{BFT2014} N. Bruin, E.~V. Flynn, and D. Testa, \emph{Descent via (3,3)-isogeny
on Jacobians of genus 2 curves}, Acta Arith.\ \textbf{165} (2014), 201--223.
\bibitem{BLP2009} N. Bernard, F. Lepr\'evost, and M. Pohst, \emph{Jacobians of
genus-2 curves with a rational point of order 11}, Experiment.\ Math.\ \textbf{18}
(2009), 65--70.
\bibitem{BoxallGrant2000} J. Boxall and D. Grant, \emph{Examples of torsion
points on genus two curves}, Trans.\ Amer.\ Math.\ Soc.\ \textbf{352} (2000),
4533--4555.
\bibitem{BSSVY} A.~R. Booker, J. Sijsling, A.~V. Sutherland, J. Voight, and D. Yasaki,
\emph{A database of genus-2 curves over the rational numbers},
LMS J.~Comput.\ Math.\ \textbf{19(A)} (2016), 235--254.
\bibitem{BookerSutherland} A.~R. Booker and A.~V. Sutherland,
\emph{Genus 2 curves of small conductor}, in preparation.
\bibitem{AC2026} J. Alessandr\`i and N. Coppola, \emph{Torsion points on
$\mathrm{GL}_2$-type abelian varieties}, arXiv:2602.21047 (v3, April 2026);
data repository \url{https://github.com/NirvanaC93/Torsion-of-GL2-abelian-varieties}.
\bibitem{CEHJMPV} E. Costa, N.~D. Elkies, S. Hashimoto, R. Jha, K. Martin,
B. Poonen, and J. Voight, \emph{Modular abelian surfaces} (data repository),
\url{https://github.com/edgarcosta/ModularAbelianSurfaces}, 2024.
% TODO(collaborators): confirm the author list of the dataset/paper before
% publication; taken from the repository (commit 687c10d).
\bibitem{Elkies2002} N.~D. Elkies, \emph{Curves of genus 2 over $\mathbb{Q}$ whose
Jacobians are absolutely simple abelian surfaces with torsion points of high order},
\url{https://people.math.harvard.edu/~elkies/g2_tors.html}, 2001--2002, updated 2010.
\bibitem{Elkies2024} N.~D. Elkies, \emph{Families of genus-2 curves with 5-torsion},
LuCaNT: LMFDB, Computation, and Number Theory, Contemp.\ Math.\ \textbf{796},
Amer.\ Math.\ Soc., 2024, 165--185.
\bibitem{Flynn1990} E.~V. Flynn, \emph{Large rational torsion on abelian varieties},
J.~Number Theory \textbf{36} (1990), 257--265.
\bibitem{Flynn1991} E.~V. Flynn, \emph{Sequences of rational torsions on abelian
varieties}, Invent.\ Math.\ \textbf{106} (1991), 433--442.
\bibitem{HLP1996} E.~W. Howe, F. Lepr\'evost, and B. Poonen, \emph{Sous-groupes de
torsion d'ordres \'elev\'es de Jacobiennes d\'ecomposables de courbes de genre 2},
C.~R.\ Acad.\ Sci.\ Paris S\'er.\ I \textbf{323} (1996), 1031--1034.
\bibitem{HLP2000} E.~W. Howe, F. Lepr\'evost, and B. Poonen, \emph{Large torsion
subgroups of split Jacobians of curves of genus two or three}, Forum Math.\
\textbf{12} (2000), 315--364.
\bibitem{Howe2015} E.~W. Howe, \emph{Genus-2 Jacobians with torsion points of
large order}, Bull.\ London Math.\ Soc.\ \textbf{47} (2015), 127--135.
\bibitem{Leprevost1991a} F. Lepr\'evost, \emph{Famille de courbes de genre 2 munies
d'une classe de diviseurs rationnels d'ordre 13},
C.~R.\ Acad.\ Sci.\ Paris S\'er.\ I Math.\ \textbf{313} (1991), 451--454.
\bibitem{Leprevost1991b} F. Lepr\'evost, \emph{Familles de courbes de genre 2 munies
d'une classe de diviseurs rationnels d'ordre 15, 17, 19 ou 21},
C.~R.\ Acad.\ Sci.\ Paris S\'er.\ I Math.\ \textbf{313} (1991), 771--774.
\bibitem{Leprevost1993} F. Lepr\'evost, \emph{Points rationnels de torsion de
jacobiennes de certaines courbes de genre 2}, C.~R.\ Acad.\ Sci.\ Paris S\'er.\ I
\textbf{316} (1993), 819--821.
\bibitem{Leprevost1995} F. Lepr\'evost, \emph{Jacobiennes de certaines courbes de
genre 2: torsion et simplicit\'e}, J.~Th\'eor.\ Nombres Bordeaux \textbf{7} (1995),
283--306.
\bibitem{Leprevost1997} F. Lepr\'evost, \emph{Sur certains sous-groupes de torsion
de jacobiennes de courbes hyperelliptiques de genre $g\ge 1$},
Manuscripta Math.\ \textbf{92} (1997), 47--63.
\bibitem{LPS2004} F. Lepr\'evost, M. Pohst, and A. Sch\"opp, \emph{Rational torsion
of $J_0(N)$ for hyperelliptic modular curves and families of Jacobians of genus 2 and
genus 3 curves with a rational point of order 5, 7 or 10}, Abh.\ Math.\ Sem.\ Univ.\
Hamburg \textbf{74} (2004), 193--203.
\bibitem{Ogg1973} A. Ogg, \emph{Rational points on certain elliptic modular curves},
Proc.\ Sympos.\ Pure Math.\ XXIV, Amer.\ Math.\ Soc., 1973, 221--231.
\bibitem{PP2015} V.~P. Platonov and M.~M. Petrunin, \emph{New curves of genus 2 over
the field of rational numbers whose Jacobians contain torsion points of high order},
Dokl.\ Math.\ \textbf{91} (2015), 220--221.
\bibitem{Nicholls2018} C. Nicholls, \emph{Descent methods and torsion on Jacobians
of higher genus curves}, Ph.D.\ thesis, University of Oxford, 2018.
\bibitem{Ogawa1994} H. Ogawa, \emph{Curves of genus 2 with a rational torsion
divisor of order 23}, Proc.\ Japan Acad.\ Ser.\ A \textbf{70} (1994), 295--298.
\bibitem{PP2012} V.~P. Platonov and M.~M. Petrunin, \emph{New orders of torsion
points in Jacobians of curves of genus 2 over the rational number field},
Dokl.\ Math.\ \textbf{85} (2012), 286--288.
\bibitem{PZP2013} V.~P. Platonov, V.~S. Zhgun, and M.~M. Petrunin,
\emph{On the simplicity of Jacobians for hyperelliptic curves of genus 2 over the
field of rational numbers with torsion points of high order},
Dokl.\ Math.\ \textbf{87} (2013), 318--321.
\bibitem{CLV2021} E. Costa, D. Lombardo, and J. Voight, \emph{Identifying central
endomorphisms of an abelian variety via Frobenius endomorphisms},
Res.\ Number Theory \textbf{7} (2021), no.~46.
\bibitem{Zywina2022} D. Zywina, \emph{Determining monodromy groups of abelian
varieties}, Res.\ Number Theory \textbf{8} (2022), no.~89.
\bibitem{LSSV2024} J. Laga, C. Schembri, A. Shnidman, and J. Voight,
\emph{Rational torsion points on abelian surfaces with quaternionic
multiplication}, Forum Math.\ Sigma \textbf{12} (2024), e92.
\end{thebibliography}
\end{document}
